% ============================================================
% CAS/Python ενότητες — Κεφάλαιο 14: Δυναμοσειρές
%
% QR code: qr_7d8143d735df.png  (→ latex/qrcodes/)
% URL: https://nikmatz.github.io/ELADIC_GR/code/ch14/
% ============================================================

\casactivbox%
  {Maxima: Δυναμοσειρές — Ακτίνα Σύγκλισης, Taylor, Πράξεις}%
  {%
    \label{cas:ch14_power_series}
    \begin{enumerate}[label=(\alph*),itemsep=2pt]
      \item \textbf{Ακτίνα σύγκλισης} $R = \lim\left|\dfrac{c_n}{c_{n+1}}\right|$:
        υπολογίστε για $\sum\frac{x^n}{n}$, $\sum\frac{x^n}{n!}$,
        $\sum n\,x^n$, $\sum\frac{x^n}{2^n}$, $\sum n!\,x^n$.
      \item \textbf{Αναπτύγματα Maclaurin} με \texttt{taylor(f,x,0,8)}:
        $e^x$, $\sin x$, $\cos x$, $\ln(1+x)$, $\dfrac{1}{1-x}$,
        $\arctan x$, $\sqrt{1+x}$. Taylor γύρω από $x_0=1$ για $e^x$.
      \item \textbf{Πράξεις}: παράγωγος $\frac{d}{dx}\!\left[\sum x^n\right]$
        $=\sum nx^{n-1}=\frac{1}{(1-x)^2}$· ολοκλήρωμα
        $\int\!\frac{1}{1+x}\,dx = \sum\frac{(-1)^n x^{n+1}}{n+1}=\ln(1+x)$.
      \item \textbf{Υπολογισμός $\pi$}: σειρά Gregory-Leibniz
        $\frac{\pi}{4}=\sum\frac{(-1)^n}{2n+1}$ για $N=10,100,1000$·
        τύπος Machin $\frac{\pi}{4}=4\arctan\frac{1}{5}-\arctan\frac{1}{239}$.
      \item Προσέγγιση $\int_0^1\!\sin(x^2)\,dx$ με δυναμοσειρά.
    \end{enumerate}
    \smallskip
    \textbf{Εντολές Maxima:}
    \texttt{taylor(f,x,x0,n)}, \texttt{diff(T,x)},
    \texttt{integrate(T,x)}, \texttt{limit()}, \texttt{sum()}
  }%
  {https://nikmatz.github.io/ELADIC_GR/code/ch14/}%
  {qr_7d8143d735df.png}

\subsection*{Δραστηριότητα Python}

\begin{pybox}{Python: Δυναμοσειρές — SymPy, Σύγκλιση, Γραφικά}%
             {https://nikmatz.github.io/ELADIC_GR/code/ch14/}%
             {qr_7d8143d735df.png}
\label{py:ch14_power_series}
\begin{enumerate}[label=(\alph*),itemsep=2pt]
  \item Πίνακας ακτίνων σύγκλισης $R$ για 6 δυναμοσειρές.
    Αναπτύγματα Maclaurin με \texttt{sympy.series(f,x,0,8)}.
  \item Γραφικά Maclaurin: σχεδιάστε $e^x$, $\sin x$, $\ln(1+x)$,
    $\arctan x$ μαζί με προσεγγίσεις τάξης $n=1,3,5,9$ —
    παρατηρήστε πού «ξεφεύγει» κάθε πολυώνυμο από την καμπύλη.
  \item Υπολογισμός $\pi$: σύγκλιση Gregory-Leibniz (αργή)
    vs τύπος Machin (ταχεία). Σφάλμα ως συνάρτηση $N$ σε log-scale.
  \item Δυναμοσειρά για ολοκλήρωμα: $\int_0^1\!\sin(x^2)\,dx$
    — σύγκριση σειράς vs \texttt{scipy.quad}. Σφάλμα ανά τάξη.
\end{enumerate}
\medskip
\textbf{Βιβλιοθήκες / Εντολές:}
\texttt{sympy.series()}, \texttt{sympy.diff()},
\texttt{sympy.integrate()}, \texttt{math.factorial()},
\texttt{scipy.integrate.quad()}, \texttt{matplotlib.pyplot}
\end{pybox}
